\subsection{热力学第二定律}
\label{sec:app-b-second-law}
\renewcommand{\thestatement}{B.3.\arabic{statement}}
\renewcommand{\theHstatement}{appB.sec03.\arabic{statement}}
\setcounter{statement}{0}

\subsubsection{熵}
\label{subsec:app-b-entropy}

热力学第二定律断言存在一个称为熵的新状态变量 $S$, 使得对于所考虑封闭系统的任意准静态变换, 都有
\begin{equation}
 \frac{\Delta Q}{T}\leq\Delta S,
 \label{eq:app-b-clausius-inequality}
\end{equation}
并且等号成立当且仅当该变换可逆. 因此, 对于一次元可逆变换, 有
\[
 \Delta Q=T\,dS.
\]
这里利用熵是状态函数这一事实, 把记号 $\Delta S$ 换成了微分记号 $dS$.

这个称为 Clausius 定理的陈述, 等价于说从冷系统向热系统传热是不允许的.

现在把热力学的这两个定律结合起来, 可以证明, 在一次元可逆变换中, 封闭系统的内能按照如下方程变化:
\begin{equation}
 dE=T\,dS-P\,dV.
 \label{eq:app-b-reversible-internal-energy}
\end{equation}

\begin{Remark}
\label{rem:app-b-irreversible-transformations}
上述结果只适用于可逆变换. 应当怎样处理不可逆变换? 可以利用能量变化只依赖于初态和终态这一事实, 尝试构造一个把系统从第一个状态带到第二个状态的可逆变换. 沿着这样一条路径, 上述关系成立, 因而可以对它积分.
\end{Remark}

\subsubsection{内能的计算}
\label{subsec:app-b-internal-energy-calculation}

现在可以给出内能关于其他变量 $T,P,V,S$ 的表达式. 为此, 只需利用每个变量的强度性质或广延性质.

考虑一个大小为 $1$ 的系统, 然后考虑同一系统变为原来的 $(1+\varepsilon)$ 倍, 同时保持其所有比例不变. 在这个系统中, 温度和压力不变, 而能量, 熵和体积都乘以 $(1+\varepsilon)$. 应用\eqref{eq:app-b-reversible-internal-energy}, 并利用每个状态中的 $T$ 和 $P$ 相同, 可得
\[
\begin{aligned}
 \varepsilon E&=(1+\varepsilon)E-E=\Delta E=T(\Delta S)-P(\Delta V)\\
 &=T((1+\varepsilon)S-S)-P((1+\varepsilon)V-V)
 =\varepsilon TS-\varepsilon PV;
\end{aligned}
\]
于是得到表达式
\[
 E=TS-PV.
\]
对此式取微分并与\eqref{eq:app-b-reversible-internal-energy}比较, 得到 Gibbs--Duhem 关系
\begin{equation}
 S\,dT-V\,dP=0.
 \label{eq:app-b-gibbs-duhem}
\end{equation}
